% ==============================================================================
% pretextual/resumo.tex
% Resumo em língua vernácula
% ==============================================================================

\setlength{\absparsep}{18pt}
\begin{resumo}
Este documento técnico apresenta a arquitetura, a especificação de engenharia, a implementação completa e a validação automatizada de uma pipeline híbrida de aquisição, análise de dados e atuação para automação industrial e Internet das Coisas (IoT), desenvolvida no âmbito de uma Situação de Aprendizagem (SA) no Centroweg Senai. Concebido inicialmente como um servidor de tempo \textit{Network Time Protocol} (NTP) baseado em microcontrolador ESP32 e receptor \textit{Global Navigation Satellite System} (GNSS), o projeto teve seu escopo integralmente redefinido após limitações práticas de recepção satelital em ambiente interno de laboratório (\textit{indoor}) e restrições de infraestrutura inviabilizarem o protótipo original, validando-se junto à coordenação a transição para uma bancada de Internet Industrial das Coisas (IIoT) autocontida e reproduzível. Na camada de campo (24\,Vcc PNP), documentada com esquemáticos elétricos em KiCad~8 e CircuiTikZ, quatro botoeiras (\texttt{B1} a \texttt{B4}) simulam sensores de presença de peça, nível crítico, sobretemperatura e parada de emergência normalmente fechada (\textit{fail-safe}), enquanto quatro LEDs (\texttt{L1} a \texttt{L4}) com resistores limitadores de 2,2\,k$\Omega$ representam motor, válvula de alívio, alarme e indicação de operação normal. Na camada de controle, programada de forma modular em \textit{Structured Text} (IEC~61131-3) no CODESYS~V3.5 (\texttt{GVL\_Modbus}, \texttt{FB\_\allowbreak Debounce\allowbreak Counter}, \texttt{FB\_\allowbreak Watchdog\allowbreak Heartbeat}, \texttt{FB\_\allowbreak Safety\allowbreak Interlock} e \texttt{PLC\_PRG} com ciclo de varredura de 10\,ms e exportação PLCopen TC6 XML v2.01), o CLP WEG PLC500 opera como servidor Modbus TCP/IP na porta 502 com filtro de \textit{debounce} de 50\,ms, contadores de eventos, \textit{watchdog} por \textit{heartbeat} e intertravamento determinístico local. Na camada de orquestração e análise, uma aplicação Node-RED~v4 estruturada em cinco abas utiliza o pacote \texttt{node-red-\allowbreak contrib-\allowbreak modbus} para realizar \textit{polling} a 500\,ms (2\,Hz) de 23 registradores de base zero, persistir séries temporais e auditorias no PostgreSQL~16 (quatro tabelas e três visões operacionais) e computar atributos estatísticos em janelas deslizantes de 60\,s com passo (\textit{stride}) de 15\,s. Esses vetores alimentam um microsserviço de aprendizado de máquina em FastAPI e \textit{scikit-learn} --- composto por detector de anomalias \texttt{Isolation\allowbreak Forest} e classificador supervisionado \texttt{HistGradient\allowbreak Boosting\allowbreak Classifier} treinados sobre 18\,h de telemetria a 2\,Hz (129.600 amostras brutas e 3.688 janelas com lacunas de purga $\ge 60$\,s) --- acoplado a um agente baseado em Modelo de Linguagem Grande (LLM) restrito por \textit{JSON Schema} Draft 2020-12 e guardião de monotonicidade de segurança. Por fim, a solução é verificada integralmente em malha fechada (\textit{Software-in-the-Loop}) nos dez casos de teste (\texttt{CT-01} a \texttt{CT-10}) por meio de um gêmeo digital Modbus TCP de 10\,ms (\texttt{plc\_\allowbreak simulator.py}).

\textbf{Palavras-chave}: Internet Industrial das Coisas. WEG PLC500. Modbus TCP/IP. Node-RED. Aprendizado de máquina. Gêmeo digital.
\end{resumo}
